\documentclass[12pt]{article}

% Source: Tutorial-8.pdf
% Style reference: MH5200_ProblemSheet2.tex

\usepackage[margin=1in]{geometry}
\usepackage{amsmath,amssymb,mathtools}
\usepackage{enumitem}
\usepackage{bm}
\usepackage{hyperref}
\usepackage{fancyhdr}
\usepackage{draftwatermark}
\usepackage{xcolor}

%------------- TITLE AND METADATA -------------

\newcommand{\CourseCode}{MH1101 }
\newcommand{\CourseName}{Calculus II}
\newcommand{\DocType}{Tutorial 8 (Week 9) -- Problems \& Solutions}

\title{
\CourseCode \CourseName \\[0.3em]
\large \DocType
}
\author{
  Academic Year 2025/2026, Semester 2 \\[0.25em]
  \textit{Quantitative Research Society @NTU}
}
\date{\today}

%------------- HEADER & FOOTER (for page 2 onward) -------------
\fancypagestyle{main}{
  \fancyhf{}
  \fancyhead[L]{\CourseCode \CourseName}
  \fancyhead[R]{\href{https://qrsntu.org}{QRS@NTU}}
  \fancyfoot[C]{\thepage}
  \renewcommand{\headrulewidth}{0.4pt}
  \renewcommand{\footrulewidth}{0pt}
}

%------------- SIMPLE FOOTER (page 1 only) -------------
\fancypagestyle{plainfirst}{
  \fancyhf{}
  \renewcommand{\headrulewidth}{0pt}
  \renewcommand{\footrulewidth}{0pt}
  \fancyfoot[C]{\scriptsize\textcolor{gray}{\textit{For academic reference only. This document is provided for educational purposes and does not constitute an official question paper, answer key, or any formally authorised academic material. Its content is not guaranteed to be complete or accurate.}}}
}

%------------- WATERMARK -------------
\SetWatermarkText{Unofficial — QRS@NTU — qrsntu.org}
\SetWatermarkScale{1}
\SetWatermarkLightness{0.99}

\begin{document}

%------------- COVER PAGE -------------
\maketitle
\thispagestyle{plainfirst}
\pagestyle{main}

\bigskip

%====================================================
% OVERVIEW
%====================================================

\begin{center}
\rule{0.8\textwidth}{0.4pt}\\[4pt]
\textbf{Overview}\\[2pt]
\rule{0.8\textwidth}{0.4pt}
\end{center}

\noindent
This tutorial develops convergence techniques for sequences (especially recursive ones) and applies standard tests to determine convergence/divergence of infinite series.

\begin{itemize}[leftmargin=*]
\item Nested radicals and rewriting recurrences to identify the limit.
\item Monotone bounded sequences from fixed-point recurrences.
\item Proving monotonicity of \(e_n=\left(1+\frac{1}{n}\right)^n\).
\item Series tests: geometric, comparison, telescoping, and divergence tests.
\end{itemize}

\bigskip

\newpage

%====================================================
% QUESTION 1
%====================================================

\section*{Question 1 (Nested radical sequence)}

\subsection*{Problem}
Find the limit of the sequence:
\[
\sqrt{2},\ \sqrt{2\sqrt{2}},\ \sqrt{2\sqrt{2\sqrt{2}}},\ \ldots
\]

\subsection*{Solution}

\subsubsection*{Method 1: Monotone + bounded, then solve the fixed-point equation}

Define \(a_1=\sqrt{2}\) and for \(n\ge 1\),
\[
a_{n+1}=\sqrt{2a_n}.
\]
First show \(0<a_n<2\) for all \(n\). For \(n=1\), \(a_1=\sqrt{2}<2\). If \(0<a_n<2\), then
\[
0<a_{n+1}=\sqrt{2a_n}<\sqrt{2\cdot 2}=2.
\]
Hence \(0<a_n<2\) for all \(n\).

Next show \(\{a_n\}\) is increasing. Since \(a_n>0\),
\[
a_{n+1}\ge a_n
\iff \sqrt{2a_n}\ge a_n
\iff 2a_n \ge a_n^2
\iff a_n(2-a_n)\ge 0,
\]
which holds because \(0<a_n<2\). Thus \(a_{n+1}\ge a_n\).

So \(\{a_n\}\) is increasing and bounded above by \(2\), hence convergent. Let \(\lim_{n\to\infty} a_n=L\). Taking limits in \(a_{n+1}=\sqrt{2a_n}\) (continuity of \(\sqrt{\cdot}\) on \((0,\infty)\)) gives
\[
L=\sqrt{2L}\quad\Rightarrow\quad L^2=2L\quad\Rightarrow\quad L(L-2)=0.
\]
Since all \(a_n>0\), we must have \(L>0\), hence \(L=2\). Therefore
\[
\boxed{\lim_{n\to\infty} a_n = 2.}
\]

\subsubsection*{Method 2: Trigonometric closed form}

Claim:
\[
a_n = 2\cos\left(\frac{\pi}{2^{n+1}}\right)\qquad (n\ge 1).
\]
For \(n=1\), \(2\cos(\pi/4)=\sqrt{2}=a_1\).

Assume \(a_n = 2\cos(\theta)\) where \(\theta=\frac{\pi}{2^{n+1}}\). Then
\[
a_{n+1}=\sqrt{2a_n}=\sqrt{4\cos\theta}=2\sqrt{\cos\theta}.
\]
Using \(\cos\frac{\theta}{2}=\sqrt{\frac{1+\cos\theta}{2}}\), and the identity \(a_{n+1}=\sqrt{2a_n}\) corresponds to the half-angle relation for cosine, one obtains
\[
2\cos\left(\frac{\theta}{2}\right)=\sqrt{2\cdot 2\cos\theta}=\sqrt{2a_n}=a_{n+1}.
\]
Thus
\[
a_{n+1}=2\cos\left(\frac{\pi}{2^{n+2}}\right),
\]
so the formula holds by induction. Hence
\[
\lim_{n\to\infty} a_n
= \lim_{n\to\infty} 2\cos\left(\frac{\pi}{2^{n+1}}\right)
=2\cos(0)=2,
\]
so \(\boxed{\lim a_n=2}\).

\newpage

%====================================================
% QUESTION 2
%====================================================

\section*{Question 2 (Monotone bounded recursion I)}

\subsection*{Problem}
Show that the sequence defined by
\[
a_1 = 1,\qquad a_{n+1} = 3 - \frac{1}{a_n}
\]
is increasing and \(0<a_n<3\) for all \(n\). Deduce that \(\{a_n\}\) is convergent and find its limit.

\subsection*{Solution}

\subsubsection*{Method 1: Invariant interval + monotonicity, then fixed point}

Let \(f(x)=3-\frac{1}{x}\) on \(x>0\), so \(a_{n+1}=f(a_n)\).

\paragraph{Step 1: Show \(0<a_n<3\) for all \(n\).}
For \(n=1\), \(a_1=1\in(0,3)\). Assume \(a_n\in(0,3)\). Then \(a_n>0\) implies \(1/a_n>0\), so \(a_{n+1}=3-\frac{1}{a_n}<3\). Also \(a_n<3\) implies \(1/a_n>1/3\), so
\[
a_{n+1}=3-\frac{1}{a_n} > 3-\frac{1}{(0^+)}\quad\text{(not useful),}
\]
but since \(a_n\ge a_1=1\) will be shown below, we can first establish positivity directly:
because \(a_n\in(0,3)\) implies \(1/a_n>1/3\), hence \(a_{n+1}=3-\frac{1}{a_n}>3-\infty\) is not immediate. Instead, note that for \(a_n\in(0,3)\),
\[
a_{n+1}>0 \iff 3-\frac{1}{a_n}>0 \iff a_n>\frac{1}{3},
\]
and indeed \(a_1=1>\frac{1}{3}\). We next show \(a_n\ge 1\) for all \(n\), which implies \(a_n>\frac{1}{3}\) and hence \(a_{n+1}>0\).

\paragraph{Step 2: Show \(\{a_n\}\) is increasing.}
Compute
\[
a_{n+1}-a_n = 3-\frac{1}{a_n}-a_n = \frac{-a_n^2+3a_n-1}{a_n}.
\]
Since \(a_n>0\), the sign is the sign of \(-a_n^2+3a_n-1\), i.e. of
\[
g(x)=-x^2+3x-1 = -(x-\alpha)(x-\beta),
\]
where
\[
\alpha=\frac{3-\sqrt{5}}{2},\qquad \beta=\frac{3+\sqrt{5}}{2}.
\]
Thus \(g(x)\ge 0\) exactly when \(x\in[\alpha,\beta]\).

Now \(a_1=1\) and \(\alpha<1<\beta\). Also \(f\) is increasing on \(x>0\) since \(f'(x)=\frac{1}{x^2}>0\). One checks that \(f([\alpha,\beta])\subseteq [\alpha,\beta]\) because \(f(\alpha)=\alpha\) and \(f(\beta)=\beta\) (they are fixed points), and \(f\) is increasing. Hence by induction \(a_n\in[\alpha,\beta]\) for all \(n\), so \(g(a_n)\ge 0\) for all \(n\), giving \(a_{n+1}\ge a_n\). In particular \(a_n\ge a_1=1\), so all \(a_n>0\), and also \(a_n\le \beta<3\). Therefore \(0<a_n<3\) for all \(n\) and \(\{a_n\}\) is increasing.

\paragraph{Step 3: Conclude convergence and compute the limit.}
Since \(\{a_n\}\) is increasing and bounded above (e.g. by \(\beta\)), it converges to some \(L\). Taking limits in \(a_{n+1}=3-\frac{1}{a_n}\) yields
\[
L = 3-\frac{1}{L}
\quad\Longleftrightarrow\quad
L^2-3L+1=0
\quad\Longleftrightarrow\quad
L\in\left\{\frac{3-\sqrt{5}}{2},\ \frac{3+\sqrt{5}}{2}\right\}.
\]
Since \(a_n\ge 1\) for all \(n\), the limit must satisfy \(L\ge 1\), hence
\[
\boxed{L=\frac{3+\sqrt{5}}{2}.}
\]

\subsubsection*{Method 2: Fixed-point iteration picture + subsequence trapping}

The map \(f(x)=3-\frac{1}{x}\) is increasing on \((0,\infty)\) and has exactly two fixed points \(\alpha<\beta\) given above. Starting from \(a_1=1\in(\alpha,\beta)\), monotonicity implies
\[
\alpha < a_1 < a_2=f(a_1) < f(\beta)=\beta,
\]
and inductively \(\alpha<a_n<\beta\) with \(a_{n+1}=f(a_n)\ge a_n\). This traps the sequence in \((\alpha,\beta)\) and forces convergence to the only fixed point in \([\;1,\infty)\), namely \(\beta=\frac{3+\sqrt{5}}{2}\). Thus \(\boxed{\lim a_n=\frac{3+\sqrt{5}}{2}}\).

\newpage

%====================================================
% QUESTION 3
%====================================================

\section*{Question 3 (Monotone bounded recursion II)}

\subsection*{Problem}
Show that the sequence defined by
\[
a_1 = 2,\qquad a_{n+1} = \frac{1}{3-a_n}
\]
satisfies \(0<a_n\le 2\), and is decreasing. Deduce that \(\{a_n\}\) is convergent and find its limit.

\subsection*{Solution}

\subsubsection*{Method 1: Invariant interval + monotonicity, then fixed point}

Let \(f(x)=\frac{1}{3-x}\), defined for \(x\neq 3\), so \(a_{n+1}=f(a_n)\).

\paragraph{Step 1: Show \(0<a_n\le 2\) for all \(n\).}
For \(n=1\), \(a_1=2\). Suppose \(0<a_n\le 2\). Then \(3-a_n\in[1,3)\), hence
\[
a_{n+1}=\frac{1}{3-a_n}\in\left( \frac{1}{3}, 1\right]\subset (0,2].
\]
So \(0<a_{n+1}\le 2\). By induction, \(0<a_n\le 2\) for all \(n\). In particular \(3-a_n>0\), so the recurrence is well-defined.

\paragraph{Step 2: Show \(\{a_n\}\) is decreasing.}
Compute
\[
a_{n+1}\le a_n
\iff \frac{1}{3-a_n}\le a_n
\iff 1\le a_n(3-a_n)
\iff a_n^2-3a_n+1\le 0.
\]
But \(a_n^2-3a_n+1=(a_n-\alpha)(a_n-\beta)\), with \(\alpha=\frac{3-\sqrt{5}}{2}\) and \(\beta=\frac{3+\sqrt{5}}{2}\).
The inequality \((a_n-\alpha)(a_n-\beta)\le 0\) holds exactly when \(a_n\in[\alpha,\beta]\).

We already have \(a_1=2\in[\alpha,\beta]\), and note that \(f\) is increasing on \((-\infty,3)\) because \(f'(x)=\frac{1}{(3-x)^2}>0\). Also \(\alpha,\beta\) are fixed points of \(f\) (they solve \(L=\frac{1}{3-L}\)). Hence \(f([\alpha,\beta])\subseteq[\alpha,\beta]\), and by induction \(a_n\in[\alpha,\beta]\). Therefore \(a_{n+1}\le a_n\) for all \(n\): the sequence is decreasing.

\paragraph{Step 3: Conclude convergence and compute the limit.}
Since \(\{a_n\}\) is decreasing and bounded below by \(0\), it converges to some \(L\ge 0\).
Taking limits in \(a_{n+1}=\frac{1}{3-a_n}\) yields
\[
L=\frac{1}{3-L}\quad\Longleftrightarrow\quad L^2-3L+1=0
\quad\Longleftrightarrow\quad
L\in\left\{\alpha,\beta\right\}.
\]
But \(a_2=\frac{1}{3-a_1}=1\), so the sequence is decreasing from \(2\) downwards and thus \(L\le 1\). Therefore \(L\neq \beta\) (since \(\beta>2\)), and we must have
\[
\boxed{L=\frac{3-\sqrt{5}}{2}.}
\]

\subsubsection*{Method 2: Two-sided squeezing using the fixed point}

Let \(L=\frac{3-\sqrt{5}}{2}\), which satisfies \(L=\frac{1}{3-L}\). Consider \(b_n=a_n-L\). Using the recurrence,
\[
b_{n+1}=a_{n+1}-L=\frac{1}{3-a_n}-\frac{1}{3-L}
=\frac{a_n-L}{(3-a_n)(3-L)}
=\frac{b_n}{(3-a_n)(3-L)}.
\]
Since \(0<a_n\le 2\), we have \(1\le 3-a_n<3\), so \((3-a_n)(3-L)\ge 1\cdot (3-L)>1\). Hence
\[
0<b_{n+1}\le \frac{1}{3-L}\,b_n=b_n\cdot L < b_n,
\]
so \(b_n\downarrow 0\), i.e. \(a_n\downarrow L\). Thus \(\boxed{\lim a_n=\frac{3-\sqrt{5}}{2}}\).

\newpage

%====================================================
% QUESTION 4
%====================================================

\section*{Question 4 (Monotonicity of \(\left(1+\frac{1}{n}\right)^n\))}

\subsection*{Problem}
Define
\[
e_n=\left(1+\frac{1}{n}\right)^n.
\]
Show that the sequence \(\{e_n\}_{n=1}^{\infty}\) is increasing.
\quad (Hint: Use the Binomial Theorem \((1+a)^n=\sum_{k=0}^{n}\binom{n}{k}a^k\).)

\subsection*{Solution}

\subsubsection*{Method 1: Binomial theorem lower bound for \(e_{n+1}\) and comparison}

Using the binomial theorem,
\[
\left(1+\frac{1}{n}\right)^n
=1+\binom{n}{1}\frac{1}{n}+\binom{n}{2}\frac{1}{n^2}+\cdots
=1+1+\frac{n(n-1)}{2n^2}+\cdots
>2.
\]
To compare consecutive terms, rewrite
\[
e_{n+1}=\left(1+\frac{1}{n+1}\right)^{n+1}
=\left(1+\frac{1}{n+1}\right)^n\left(1+\frac{1}{n+1}\right).
\]
Now note that \(\left(1+\frac{1}{n+1}\right) > \left(1+\frac{1}{n}\right)^{\frac{n}{n+1}}\) (this can be shown by expanding both sides with binomial theorem and comparing coefficients, or by the convexity argument in Method 2). Multiplying by \(\left(1+\frac{1}{n+1}\right)^{n/(n+1)}\) yields \(e_{n+1}>e_n\). Hence \(\{e_n\}\) is increasing.

\smallskip
\noindent
\textit{Remark.} A fully algebraic (binomial-only) coefficient comparison is possible but tends to be longer; Method 2 is typically the cleanest rigorous route.

\subsubsection*{Method 2: Calculus on a continuous extension}

Define a function for real \(x>0\):
\[
\phi(x)=x\ln\left(1+\frac{1}{x}\right).
\]
Then \(e_n=\exp(\phi(n))\). It suffices to show \(\phi(x)\) is increasing for \(x>0\).

Differentiate:
\begin{align*}
\phi'(x)
&=\ln\left(1+\frac{1}{x}\right)+x\cdot \frac{1}{1+\frac{1}{x}}\cdot\left(-\frac{1}{x^2}\right)\\
&=\ln\left(1+\frac{1}{x}\right)-\frac{1}{x+1}.
\end{align*}
Use the inequality \(\ln(1+u)\ge \frac{u}{1+u}\) for \(u>0\) (e.g. by concavity of \(\ln\), or by considering \(h(u)=\ln(1+u)-\frac{u}{1+u}\) and checking \(h'(u)\ge 0\)). With \(u=\frac{1}{x}\),
\[
\ln\left(1+\frac{1}{x}\right)\ge \frac{\frac{1}{x}}{1+\frac{1}{x}}=\frac{1}{x+1}.
\]
Therefore \(\phi'(x)\ge 0\) for \(x>0\), so \(\phi\) is increasing, hence \(e_n=\exp(\phi(n))\) is increasing:
\[
\boxed{e_{n+1}>e_n\ \text{for all }n\ge 1.}
\]

\newpage

%====================================================
% QUESTION 5
%====================================================

\section*{Question 5 (Series: convergence/divergence and sums)}

\subsection*{Problem}
Determine whether the series is convergent or divergent. If it is convergent, find its sum.
\begin{enumerate}[label=(\roman*)]
\item \(\displaystyle \sum_{n=1}^{\infty}\frac{(-3)^{\,n-1}}{4^{n}}\).
\item \(\displaystyle \sum_{n=1}^{\infty}\frac{6\cdot 2^{2n-1}}{3^{n}}\).
\item \(\displaystyle \sum_{n=1}^{\infty}\frac{e^{n}}{n^{2}}\).
\item \(\displaystyle \sum_{n=2}^{\infty}\left(\frac{1}{e^{n}}+\frac{2}{n^{2}-1}\right)\).
\item \(\displaystyle \sum_{n=1}^{\infty}\frac{n^{3}+n^{2}}{n^{3}-2n+5}\).
\item \(\displaystyle \sum_{n=1}^{\infty}\left(\frac{3}{5}n+\frac{2}{n}\right)^{2}\).
\end{enumerate}

\subsection*{Solution}

\subsubsection*{Method 1: Standard tests (geometric, divergence test, telescoping, comparison)}

\begin{enumerate}[label=(\roman*)]
\item Rewrite as a geometric series:
\[
\sum_{n=1}^{\infty}\frac{(-3)^{n-1}}{4^{n}}
=\frac{1}{4}\sum_{n=1}^{\infty}\left(\frac{-3}{4}\right)^{n-1}.
\]
Since \(\left|\frac{-3}{4}\right|<1\),
\[
\sum_{n=1}^{\infty}\left(\frac{-3}{4}\right)^{n-1}=\frac{1}{1-\left(-\frac{3}{4}\right)}=\frac{1}{1+\frac{3}{4}}=\frac{4}{7}.
\]
Thus
\[
\boxed{\sum_{n=1}^{\infty}\frac{(-3)^{n-1}}{4^{n}}=\frac{1}{7}.}
\]

\item Simplify the general term:
\[
\frac{6\cdot 2^{2n-1}}{3^{n}}
=\frac{6}{2}\cdot \frac{4^{n}}{3^{n}}
=3\left(\frac{4}{3}\right)^{n}.
\]
Since \(\frac{4}{3}>1\), the terms do not approach \(0\) and in fact grow without bound. Hence the series diverges:
\[
\boxed{\text{Divergent.}}
\]

\item Since \(e^{n}/n^{2}\to \infty\), in particular it does not tend to \(0\). Therefore by the \(n\)-th term test the series diverges:
\[
\boxed{\text{Divergent.}}
\]

\item Split:
\[
\sum_{n=2}^{\infty}\left(\frac{1}{e^{n}}+\frac{2}{n^{2}-1}\right)
=\sum_{n=2}^{\infty}\frac{1}{e^{n}}+\sum_{n=2}^{\infty}\frac{2}{n^{2}-1}.
\]
First is geometric with ratio \(1/e\):
\[
\sum_{n=2}^{\infty}\frac{1}{e^{n}}=\frac{1/e^{2}}{1-1/e}=\frac{1}{e(e-1)}.
\]
Second telescopes:
\[
\frac{2}{n^{2}-1}=\frac{2}{(n-1)(n+1)}=\frac{1}{n-1}-\frac{1}{n+1}.
\]
Hence, for \(N\ge 2\),
\begin{align*}
\sum_{n=2}^{N}\left(\frac{1}{n-1}-\frac{1}{n+1}\right)
&=\left(1+\frac{1}{2}\right)-\left(\frac{1}{N}+\frac{1}{N+1}\right)
\to \frac{3}{2}.
\end{align*}
Therefore
\[
\boxed{\sum_{n=2}^{\infty}\left(\frac{1}{e^{n}}+\frac{2}{n^{2}-1}\right)
=\frac{1}{e(e-1)}+\frac{3}{2}.}
\]

\item Check the term limit:
\[
\frac{n^{3}+n^{2}}{n^{3}-2n+5}\to \frac{1+0}{1+0+0}=1\neq 0,
\]
so the series diverges by the \(n\)-th term test:
\[
\boxed{\text{Divergent.}}
\]

\item The general term is
\[
\left(\frac{3}{5}n+\frac{2}{n}\right)^2 \sim \left(\frac{3}{5}n\right)^2=\frac{9}{25}n^2,
\]
so it does not go to \(0\). Hence the series diverges by the \(n\)-th term test:
\[
\boxed{\text{Divergent.}}
\]
\end{enumerate}

\subsubsection*{Method 2: Alternative confirmations (ratio/root tests and partial sums)}

\begin{enumerate}[label=(\roman*)]
\item For (i), treat it directly as geometric with first term \(a=\frac14\) and ratio \(r=-\frac34\). Then \(S=\frac{a}{1-r}=\frac{1/4}{1+3/4}=\frac17\).

\item For (ii), ratio test on \(u_n=3(4/3)^n\) gives \(u_{n+1}/u_n=4/3>1\), so \(u_n\not\to 0\) and the series diverges.

\item For (iii), ratio test:
\[
\frac{u_{n+1}}{u_n}=\frac{e^{n+1}/(n+1)^2}{e^{n}/n^2}=e\left(\frac{n}{n+1}\right)^2\to e>1,
\]
so the series diverges.

\item For (iv), compute the partial sum explicitly:
\[
\sum_{n=2}^{N}\frac{2}{n^2-1}=\frac{3}{2}-\frac{1}{N}-\frac{1}{N+1}\to\frac{3}{2},
\]
and combine with the finite geometric tail for \(\sum_{n=2}^\infty e^{-n}\).

\item For (v), since terms tend to \(1\), the partial sums grow at least like \(\sum 1\), so diverge.

\item For (vi), since the terms grow like \(cn^2\), the partial sums grow at least like \(\sum cn^2\), hence diverge.
\end{enumerate}

\end{document}
